\chapter{Moment Polytopes and Topology}
\label{chap:moment-polytopes-topology}

\begin{chapteroverview}
Until now, a lattice polytope has played several algebraic and
combinatorial roles. Its lattice points indexed global sections, its
normal fan constructed a projective toric variety, and its edges and
area encoded intersection numbers on surfaces.

Over \(\CC\), the same polytope acquires a new meaning: it is the image
of a continuous map from the toric variety itself. The mechanism begins
with the compact torus
\[
K_N\simeq(S^1)^n\subseteq T_N\simeq(\CC^*)^n.
\]
The compact torus changes the phases of monomial coordinates but not
their absolute values. Normalized coordinate magnitudes therefore give
convex coefficients, and the lattice weights of the coordinates combine
to produce a point of the polytope.

This chapter develops the progression
\[
\boxed{
\begin{gathered}
\text{moment map}
\longrightarrow
\text{polytope as orbit space}
\longrightarrow
\text{cells and Betti numbers}
\\[-1mm]
\longrightarrow
\text{the cohomology ring from the fan}.
\end{gathered}}
\]
Throughout the chapter, we set \(\Bbbk=\CC\) and use the classical
topology.
\end{chapteroverview}

\begin{threedictionaries}
The topology of a projective toric variety is governed by the same
weights that defined its projective coordinates.

\medskip

\begin{center}
\small
\begin{tabularx}{\linewidth}{@{}
  >{\raggedright\arraybackslash}X
  >{\raggedright\arraybackslash}X
  >{\raggedright\arraybackslash}X@{}}
\toprule
\textbf{Algebra}
&
\textbf{Geometry and topology}
&
\textbf{Combinatorics}
\\
\midrule
monomial coordinates \(\character{m_i}\)
&
projective toric embedding
&
lattice weights \(m_i\in M\)
\\
coordinate phases
&
compact-torus orbits
&
directions that disappear under \(\mu\)
\\
normalized values \(\lvert z_i\rvert^2\)
&
moment map
&
convex coefficients
\\
torus-invariant divisors
&
degree-two cohomology classes
&
rays and Stanley--Reisner relations
\\
\bottomrule
\end{tabularx}
\end{center}
\end{threedictionaries}

\begin{surfacelaboratory}
For a projective toric surface, the moment image is a polygon. Over an
interior point sits a two-dimensional compact torus; over an edge, one
circle direction has collapsed; over a vertex, the entire compact torus
has collapsed to a fixed point.

A toric blowup cuts a corner from the moment polygon. Later in the
chapter, the same corner cut will produce one new two-dimensional
cohomology class and change the intersection form by adjoining an
exceptional class of square \(-1\).
\end{surfacelaboratory}

\section{Compact torus actions and the moment map}
\label{sec:compact-torus-moment-map}
\index{compact torus}
\index{moment map}
\index{moment polytope}

Throughout this chapter, we work over
\[
\Bbbk=\CC
\]
and use the classical topology on complex varieties. This change of
field matters: an element of \(\CC^*\) has both a positive size and a
phase, and the phase directions form a compact torus.

\subsection{The compact torus inside the algebraic torus}

Recall that
\[
T_N=\Hom_{\ZZ}(M,\CC^*).
\]
Let
\[
S^1=\{z\in\CC:\lvert z\rvert=1\}.
\]
The \emph{compact torus} associated to \(N\) is
\[
K_N:=\Hom_{\ZZ}(M,S^1)\subseteq T_N.
\]
After choosing a basis of \(M\), these groups become
\[
T_N\simeq(\CC^*)^n,
\qquad
K_N\simeq(S^1)^n.
\]

Every nonzero complex number has a unique polar decomposition
\[
z=\lvert z\rvert\frac{z}{\lvert z\rvert}.
\]
Taking the logarithm of its positive factor gives
\[
\CC^*\simeq\RR\times S^1,
\qquad
z\longmapsto
\left(\log\lvert z\rvert,\frac{z}{\lvert z\rvert}\right).
\]
Applying \(\Hom_{\ZZ}(M,-)\) yields the intrinsic toric version.

\begin{proposition}[Radial and angular directions]
\label{prop:radial-angular-directions}
There is an isomorphism of real Lie groups
\[
\boxed{T_N\simeq N_{\RR}\times K_N.}
\]
The first factor records absolute values of characters, while the second
records their phases.
\end{proposition}

\begin{proof}
For \(t\in T_N\), define
\[
\nu_t:M\longrightarrow\RR,
\qquad
m\longmapsto\log\lvert\character{m}(t)\rvert.
\]
Because absolute value turns products into products and logarithms turn
products into sums, \(\nu_t\) is a homomorphism. Hence
\[
\nu_t\in\Hom_{\ZZ}(M,\RR)=N_{\RR}.
\]
The phase
\[
m\longmapsto
\frac{\character{m}(t)}{\lvert\character{m}(t)\rvert}
\]
is an element of \(K_N\). These two assignments are inverse to the map
obtained by multiplying radial and angular characters.
\end{proof}

\begin{observation}
No background in Lie groups is needed here. The proposition is simply
the coordinatewise statement
\[
(\CC^*)^n
\simeq
(\RR_{>0})^n\times(S^1)^n
\simeq
\RR^n\times(S^1)^n.
\]
The noncompact factor \(N_{\RR}\) moves toward and away from the toric
boundary. The compact factor \(K_N\) rotates without changing any
absolute value.
\end{observation}

\begin{pauseandthink}
Let \(k\in K_N\). Why does
\[
\lvert\character{m}(k)\rvert=1
\]
for every \(m\in M\)? Explain why a formula involving only absolute
values of projective coordinates should be constant on \(K_N\)-orbits.
\end{pauseandthink}

\subsection{A projective toric variety with weighted coordinates}

Let
\[
A=\{m_0,\dots,m_r\}\subseteq M
\]
be a finite set of lattice points. As in Section~5.2,
it defines a monomial map
\[
\varphi_A:T_N\longrightarrow\PP^r,
\qquad
t\longmapsto
[\character{m_0}(t):\cdots:\character{m_r}(t)].
\]
Write
\[
X_A:=\overline{\varphi_A(T_N)}\subseteq\PP^r.
\]
The coordinates carry the weights \(m_0,\dots,m_r\). Accordingly,
\(K_N\) acts on \(\PP^r\) by
\[
k\cdot[z_0:\cdots:z_r]
=
[\character{m_0}(k)z_0:\cdots:\character{m_r}(k)z_r].
\]
Since every factor \(\character{m_i}(k)\) has absolute value one, the
action changes phases but not coordinate magnitudes.

This suggests discarding the phases and retaining only the relative
sizes
\[
\lvert z_0\rvert^2,\dots,\lvert z_r\rvert^2.
\]
Once normalized by their sum, these become coefficients of a convex
combination of the weights \(m_i\).

\begin{definition}[Projective moment map]
\label{def:projective-moment-map}
The \emph{moment map} associated to \(A\) is
\[
\mu_A:X_A\longrightarrow M_{\RR},
\qquad
\boxed{
\mu_A([z_0:\cdots:z_r])
=
\frac{\sum_{i=0}^r\lvert z_i\rvert^2m_i}
     {\sum_{i=0}^r\lvert z_i\rvert^2}.}
\]
\end{definition}

\begin{proposition}[The first properties of the moment map]
\label{prop:first-moment-map-properties}
The map \(\mu_A\) is well defined, continuous, and invariant under the
action of \(K_N\). Moreover,
\[
\mu_A(X_A)\subseteq\Conv(A).
\]
\end{proposition}

\begin{proof}
Multiplying all homogeneous coordinates by \(c\in\CC^*\) multiplies
both the numerator and denominator by \(\lvert c\rvert^2\), so the
quotient is unchanged. Thus the formula is well defined on projective
space.

For \(k\in K_N\),
\[
\lvert\character{m_i}(k)z_i\rvert^2=\lvert z_i\rvert^2,
\]
which proves invariance. Continuity is immediate. Finally, the numbers
\[
p_i=\frac{\lvert z_i\rvert^2}{\sum_j\lvert z_j\rvert^2}
\]
are nonnegative and satisfy \(\sum_i p_i=1\). Hence
\[
\mu_A([z])=\sum_i p_i m_i
\]
is a convex combination of the lattice points in \(A\).
\end{proof}

The important surprise is that the inclusion in
\cref{prop:first-moment-map-properties} is an equality. Although
\(X_A\) occupies only a small part of \(\PP^r\), it contains exactly
the right points to fill the entire convex hull of its weights.

For a projective toric variety \(X_P\), its \emph{positive real part}
\((X_P)_{\geq0}\) is the closure of the points of the dense torus on
which every character takes a positive real value. It is a concrete
slice transverse to the compact-torus orbits.

\begin{theorem}[The moment image is the polytope]
\label{thm:moment-image-is-polytope}
Let \(P\subseteq M_{\RR}\) be a full-dimensional lattice polytope, and
assume the lattice points
\[
P\cap M=\{m_0,\dots,m_r\}
\]
define a projective embedding
\[
X_P\hookrightarrow\PP^r.
\]
Then the associated moment map satisfies
\[
\boxed{\mu_P(X_P)=P.}
\]
More precisely, \(\mu_P\) maps the positive real part of \(X_P\)
homeomorphically onto \(P\).
\end{theorem}

\noindent
\textit{Reference.}
Cox--Little--Schenck, Theorem~12.2.5 proves the homeomorphism for the
algebraic moment map. Formula~(12.2.4) gives the squared-norm moment map
used here, and Exercise~12.2.8 establishes the same conclusion for this
version.

\begin{proof}[Conceptual proof sketch]
On the positive real part of the dense torus, write a point as
\(\exp(u)\) with \(u\in N_{\RR}\). Then
\[
\lvert\character{m_i}(\exp(u))\rvert^2
=
\exp(2\langle m_i,u\rangle),
\]
so
\[
\mu_P(u)
=
\frac{\sum_i e^{2\langle m_i,u\rangle}m_i}
     {\sum_i e^{2\langle m_i,u\rangle}}.
\]
This is the gradient, up to a factor of two, of the strictly convex
function
\[
F(u)=\log\left(\sum_i e^{2\langle m_i,u\rangle}\right).
\]
Strict convexity makes the gradient injective. Convex duality shows
that its image is the interior of \(P\).

As \(u\) tends to infinity in a chosen direction, the exponential
weights concentrate on the face where the corresponding linear
functional is largest. Thus boundary directions in the torus land on
faces of \(P\), and taking the closure fills all of \(P\).
\end{proof}

\begin{center}
\fbox{\parbox{0.84\linewidth}{\centering
The moment map converts \emph{relative coordinate magnitudes} into
\emph{barycentric coordinates in a lattice polytope}.}}
\end{center}

\subsection{The interval inside \texorpdfstring{\(\PP^1\)}{P1}}

\begin{example}[The moment map of \(\PP^1\)]
\label{ex:moment-map-projective-line}
Take
\[
A=\{0,1\}\subseteq\ZZ.
\]
Then \(X_A=\PP^1\), and
\[
\mu([z_0:z_1])
=
\frac{\lvert z_1\rvert^2}
     {\lvert z_0\rvert^2+\lvert z_1\rvert^2}
\in[0,1].
\]
Every \(s\in[0,1]\) occurs: one may take
\[
[z_0:z_1]=[\sqrt{1-s}:\sqrt{s}].
\]
The two endpoints come from the fixed points
\[
[1:0]
\qquad\text{and}\qquad
[0:1].
\]
All other values retain one circle of phase information.
\end{example}

\begin{figure}[ht]
\centering
\begin{tikzpicture}[scale=0.92]
  % The sphere model of P^1.
  \begin{scope}[xshift=-3.7cm]
    \node[font=\bfseries] at (0,2.2) {\(\PP^1\simeq S^2\)};
    \draw[thick] (0,0) circle (1.45);
    \draw[black!45] (-1.45,0) arc[start angle=180,end angle=360,
      x radius=1.45,y radius=0.36];
    \draw[black!30,dashed] (1.45,0) arc[start angle=0,end angle=180,
      x radius=1.45,y radius=0.36];
    \draw[black!35] (-1.19,0.82) arc[start angle=180,end angle=360,
      x radius=1.19,y radius=0.24];
    \draw[black!25,dashed] (1.19,0.82) arc[start angle=0,end angle=180,
      x radius=1.19,y radius=0.24];
    \draw[black!35] (-1.19,-0.82) arc[start angle=180,end angle=360,
      x radius=1.19,y radius=0.24];
    \draw[black!25,dashed] (1.19,-0.82) arc[start angle=0,end angle=180,
      x radius=1.19,y radius=0.24];
    \fill (0,1.45) circle (2pt);
    \fill (0,-1.45) circle (2pt);
    \node[align=center] at (0,-2.08)
      {latitude circles are\\\(S^1\)-orbits};
  \end{scope}

  % Arrow.
  \draw[-{Latex[length=3mm]},very thick,DictionaryBlue]
    (-1.35,0)--(1.35,0) node[midway,above] {\(\mu\)};

  % The moment interval.
  \begin{scope}[xshift=3.7cm]
    \node[font=\bfseries] at (0,2.2) {moment interval};
    \draw[very thick,DictionaryBlue] (0,-1.45)--(0,1.45);
    \foreach \y in {-1.45,-0.72,0,0.72,1.45}
      \fill[DictionaryBlue] (0,\y) circle (2.1pt);
    \node[right] at (0,1.45) {\(1\)};
    \node[right] at (0,-1.45) {\(0\)};
    \node[align=center] at (0,-2.08)
      {endpoints are fixed;\\interior fibers are circles};
  \end{scope}
\end{tikzpicture}
\caption{The moment map collapses each circle orbit in \(\PP^1\) to
one point of the interval. Section~8.2 will show
that this behavior is the model for every projective toric variety.}
\label{fig:moment-map-projective-line}
\end{figure}

\subsection{Triangles, rectangles, and corner cuts}

\begin{example}[The triangle of \(\PP^2\)]
\label{ex:moment-triangle-projective-plane}
For
\[
A=\{0,e_1,e_2\},
\]
the moment map is
\[
\mu([z_0:z_1:z_2])
=
\frac{\lvert z_1\rvert^2}{S}e_1
+
\frac{\lvert z_2\rvert^2}{S}e_2,
\qquad
S=\lvert z_0\rvert^2+\lvert z_1\rvert^2+\lvert z_2\rvert^2.
\]
Thus its two coordinates are nonnegative and have sum at most one.
The image is the standard triangle
\[
\Delta_2=\Conv(0,e_1,e_2).
\]
Conversely, a point \((a,b)\in\Delta_2\) is the image of
\[
[\sqrt{1-a-b}:\sqrt a:\sqrt b].
\]
So here the theorem is visible directly from barycentric coordinates.
\end{example}

\begin{example}[The rectangle of \(\PP^1\times\PP^1\)]
\label{ex:moment-rectangle-product-lines}
For the line bundle \(\cO(1,1)\), the moment map can be written without
using Segre coordinates:
\[
\mu([x_0:x_1],[y_0:y_1])
=
\left(
\frac{\lvert x_1\rvert^2}{\lvert x_0\rvert^2+\lvert x_1\rvert^2},
\frac{\lvert y_1\rvert^2}{\lvert y_0\rvert^2+\lvert y_1\rvert^2}
\right).
\]
Its image is the square \([0,1]^2\). More generally,
\(\cO(a,b)\) gives a rectangle with side lengths \(a\) and \(b\).
\end{example}

\begin{example}[A blowup becomes a corner cut]
\label{ex:moment-polygon-blowup-projective-plane}
On
\[
X=\Bl_p(\PP^2),
\]
the divisor \(2H-E\) has section polytope
\[
P_{2H-E}
=
\{(a,b)\in\RR^2:a,b\geq0,\ 1\leq a+b\leq2\}.
\]
Its five lattice points give an embedding into \(\PP^4\), and
\cref{thm:moment-image-is-polytope} says that the moment image is
exactly this truncated triangle. The exceptional divisor is visible
as the new edge created by cutting off the corner of the triangle for
\(2H\).
\end{example}

\begin{figure}[ht]
\centering
\begin{tikzpicture}[scale=0.88]
  % P^2.
  \begin{scope}[xshift=-4.6cm]
    \node[font=\bfseries] at (1.15,2.55) {\(\PP^2\)};
    \fill[DictionaryBlue!12] (0,0)--(2.1,0)--(0,2.1)--cycle;
    \draw[polytope,DictionaryBlue] (0,0)--(2.1,0)--(0,2.1)--cycle;
    \draw[->,black!35] (-0.22,0)--(2.55,0);
    \draw[->,black!35] (0,-0.22)--(0,2.55);
    \foreach \p in {(0,0),(2.1,0),(0,2.1)}
      \fill \p circle (2.1pt);
    \node at (1.05,-0.55) {triangle};
  \end{scope}

  % P^1 x P^1.
  \begin{scope}[xshift=0cm]
    \node[font=\bfseries] at (1.05,2.55) {\(\PP^1\times\PP^1\)};
    \fill[PredictionTeal!12] (0,0) rectangle (2.1,2.1);
    \draw[polytope,PredictionTeal] (0,0) rectangle (2.1,2.1);
    \draw[->,black!35] (-0.22,0)--(2.55,0);
    \draw[->,black!35] (0,-0.22)--(0,2.55);
    \foreach \p in {(0,0),(2.1,0),(0,2.1),(2.1,2.1)}
      \fill \p circle (2.1pt);
    \node at (1.05,-0.55) {rectangle};
  \end{scope}

  % Blowup.
  \begin{scope}[xshift=4.6cm]
    \node[font=\bfseries] at (1.15,2.55) {\(\Bl_p(\PP^2)\)};
    \fill[SurfaceGreen!12]
      (0,0.63)--(0.63,0)--(2.1,0)--(0,2.1)--cycle;
    \draw[polytope,SurfaceGreen]
      (0,0.63)--(0.63,0)--(2.1,0)--(0,2.1)--cycle;
    \draw[WarningRed,dashed] (0,0)--(0,0.63);
    \draw[WarningRed,dashed] (0,0)--(0.63,0);
    \draw[->,black!35] (-0.22,0)--(2.55,0);
    \draw[->,black!35] (0,-0.22)--(0,2.55);
    \foreach \p in {(0,0.63),(0.63,0),(2.1,0),(0,2.1)}
      \fill \p circle (2.1pt);
    \node[rotate=-45,SurfaceGreen,font=\small] at (0.52,0.52)
      {new edge};
    \node at (1.05,-0.55) {corner-cut polygon};
  \end{scope}
\end{tikzpicture}
\caption{Moment images of three projective toric surfaces. The same
polygons previously indexed sections and determined normal fans; over
\(\CC\), they are also the actual images of continuous maps from the
varieties.}
\label{fig:moment-images-toric-surfaces}
\end{figure}

\begin{observation}[Dependence on the chosen line bundle]
\label{obs:moment-map-depends-on-polarization}
The moment polytope belongs to a polarized toric variety, not merely to
the underlying variety. Translating every weight by \(m_0\in M\)
changes the moment map by
\[
\mu_{A+m_0}=\mu_A+m_0,
\]
while replacing a line bundle by its \(d\)-th tensor power dilates the
polytope:
\[
P\longmapsto dP.
\]
Thus the precise position and scale of the moment image depend on the
chosen projective embedding. Its normal fan, and hence the underlying
toric variety when the divisor is ample, does not change.
\end{observation}

\begin{observation}[Why symplectic geometers call this a moment map]
\label{obs:symplectic-moment-map}
The Fubini--Study form on projective space makes the compact torus
action Hamiltonian, and the formula in
\cref{def:projective-moment-map} is its symplectic moment map. We will
not need the general language of symplectic forms or Hamiltonian vector
fields. For us, the important facts are the explicit weighted-average
formula, its invariance under \(K_N\), and its polytope image.
\end{observation}

\clearpage
\subsection{What information has the map forgotten?}

Because \(\mu_P\) is constant on compact-torus orbits, it factors
through the quotient map:
\[
\begin{tikzcd}[column sep=large,row sep=large]
X_P \arrow[r,"\mu_P"] \arrow[d,two heads]
  & P \\
X_P/K_N \arrow[ur,dashed,"\overline{\mu}_P"'] & {}
\end{tikzcd}
\]
The map \(\overline{\mu}_P\) is surjective by
\cref{thm:moment-image-is-polytope}. What remains unclear is whether
it is injective. The examples suggest that the only information
forgotten by \(\mu_P\) is phase information, but the number of surviving
phases changes along the boundary of \(P\).

\begin{prediction}
The map \(\mu_P\) is constant on every \(K_N\)-orbit. In the examples,
its fibers seem to change dimension as the image point moves from the
interior of \(P\) to an edge and then to a vertex. For a point in the
relative interior of a codimension-\(k\) face, which \(k\) circle
directions should collapse? Formulate a prediction for the dimension
of the compact-torus orbit in that fiber.
\end{prediction}

\begin{guidingquestion}
Does the moment map identify \emph{exactly} the points lying in the same
compact-torus orbit? Equivalently, is the quotient space
\[
X_P/K_N
\]
actually homeomorphic to the polytope \(P\)?
\end{guidingquestion}

\section{The polytope as an orbit space}
\label{sec:polytope-orbit-space}
\index{orbit space!moment polytope}
\index{moment map!fibers}
\index{face stratum}
\index{stabilizer!of a polytope face}

The answer to the question at the end of the previous section is yes.
The moment map forgets the angular coordinates, but it makes no
additional identifications: two points have the same moment image
precisely when they lie in the same compact-torus orbit.

The shape of the fiber is not constant. Over the interior of the
polytope it is the full compact torus \(K_N\). Each time the image
reaches a new boundary direction, one or more circle directions become
stabilizers and disappear from the orbit. Our first task is to match
those disappearing directions with faces of the polytope.

\subsection{Faces and torus strata}

Let \(P\subseteq M_{\RR}\) be a full-dimensional lattice polytope whose
lattice points define a projective toric embedding. For a face
\[
Q\preceq P,
\]
write \(\sigma_Q\) for its normal cone. Recall from Section~5.4 that
\[
\dim Q+\dim\sigma_Q=n.
\]
Since the fan of \(X_P\) is the normal fan \(\Sigma_P\), the
orbit--cone correspondence attaches an algebraic torus orbit
\[
O_Q:=O(\sigma_Q)\subseteq X_P
\]
to the face \(Q\). Its dimension is
\[
\dim O_Q
=n-\dim\sigma_Q
=\dim Q.
\]
Thus the face--orbit correspondence preserves dimension, even though
the intermediate passage from faces to normal cones reverses it.

In the projective coordinates indexed by \(P\cap M\), the orbit
\(O_Q\) has a particularly useful description:
\[
x=[z_m]_{m\in P\cap M}\in O_Q
\quad\Longrightarrow\quad
z_m\neq0
\ \Longleftrightarrow\ 
m\in Q\cap M.
\]
Coordinates whose weights lie outside \(Q\) vanish, while the
coordinates indexed by lattice points of \(Q\) remain nonzero. The
closure
\[
X_Q:=\overline{O_Q}
\]
is itself the projective toric variety associated to the face \(Q\).

\begin{proposition}[The moment map respects face strata]
\label{prop:moment-map-respects-face-strata}
For every face \(Q\preceq P\),
\[
\boxed{
\mu_P(X_Q)=Q,
\qquad
\mu_P(O_Q)=\relint(Q).
}
\]
Equivalently,
\[
\mu_P^{-1}\bigl(\relint(Q)\bigr)=O_Q.
\]
\end{proposition}

\noindent
\textit{Reference.}
The face-by-face argument in Cox--Little--Schenck,
Theorem~12.2.5 proves the statement for the algebraic moment map;
Exercise~12.2.8 gives the same result for the squared-norm moment map
used here.

\begin{proof}[Conceptual proof sketch]
On \(X_Q\), precisely the coordinates with weights in \(Q\) can be
nonzero, so the moment-map formula restricts to the moment map for
\(Q\). Thus \(X_Q\) maps onto \(Q\). On the dense orbit \(O_Q\), every
one of these coordinates is nonzero, so the resulting convex
combination cannot lie in a proper face of \(Q\): equality in a
supporting inequality would force all nonzero weights into that face.
Hence \(O_Q\) maps onto \(\relint(Q)\), while its smaller
orbit strata map to the boundary faces.
\end{proof}

The proposition gives the old orbit--face dictionary a new topological
column. It also explains why the moment map knows exactly which
algebraic torus orbit contains a point: one first finds the unique face
whose relative interior contains its image.

\subsection{The compact stabilizer of a face}

The algebraic orbit \(O_Q\) is a complex torus of dimension \(\dim Q\).
Its maximal compact subgroup will be the fiber of the moment map over
\(\relint(Q)\).

To describe it intrinsically, set
\[
N_{\sigma_Q}
=N\cap\Span_{\RR}(\sigma_Q),
\qquad
N(\sigma_Q)=N/N_{\sigma_Q}.
\]
The character lattice of \(O_Q\) is
\[
\sigma_Q^\perp\cap M.
\]
Restriction of characters gives a surjective map of compact tori
\[
K_N
\longrightarrow
K_{N(\sigma_Q)}
=\Hom_{\ZZ}(\sigma_Q^\perp\cap M,S^1).
\]

\begin{definition}[Compact stabilizer of a face]
\label{def:compact-stabilizer-face}
The \emph{compact stabilizer} associated to \(Q\) is
\[
K_Q
:=
\left\{
k\in K_N:
\character{m}(k)=1
\text{ for every }m\in\sigma_Q^\perp\cap M
\right\}.
\]
Equivalently, \(K_Q\) is the kernel in the exact sequence
\[
1\longrightarrow K_Q
\longrightarrow K_N
\longrightarrow K_{N(\sigma_Q)}
\longrightarrow1.
\]
\end{definition}

The lattice \(\sigma_Q^\perp\cap M\) has rank \(\dim Q\). Therefore
\[
\boxed{
K_Q\simeq(S^1)^{\codim Q},
\qquad
K_N/K_Q\simeq(S^1)^{\dim Q}.
}
\]
Every point of \(O_Q\) has stabilizer \(K_Q\) under the \(K_N\)-action.
Indeed, the algebraic torus is abelian, so all points in one algebraic
orbit have the same stabilizer, and the action on \(O_Q\) factors
through its compact torus \(K_{N(\sigma_Q)}\).

\noindent
\textit{Reference.}
This is the compact-torus form of Cox--Little--Schenck,
Lemma~3.2.5 and Proposition~12.2.3(b).

For a facet \(F\), the normal cone is the ray generated by its
primitive inward normal \(u_F\). The subgroup \(K_F\) is the circle
\[
S^1\longrightarrow K_N,
\qquad
e^{i\theta}\longmapsto\exp(i\theta u_F).
\]
This is the precise meaning of the statement that the circle normal
to a facet collapses along that facet.

\begin{pauseandthink}
Let \(P\) be a polygon and let \(v\) be a vertex where the two edges
\(F_1\) and \(F_2\) meet. The stabilizers over the interiors of the
edges are the circles generated by \(u_{F_1}\) and \(u_{F_2}\).

What is the stabilizer over \(v\)? What dimension should the moment-map
fiber over \(v\) have?
\end{pauseandthink}

\subsection{Every fiber is one compact-torus orbit}

We can now state the central result. Its key input is the positive real
part introduced in Section~8.1. It gives one canonical point over each
point of \(P\); polar decomposition supplies all remaining points in
the fiber.

\begin{theorem}[Fibers of the moment map]
\label{thm:moment-fibers-compact-orbits}
Let \(p\in P\), and let \(Q\preceq P\) be the unique face satisfying
\[
p\in\relint(Q).
\]
Then the fiber of the moment map is a single compact-torus orbit:
\[
\boxed{
\mu_P^{-1}(p)=K_N\cdot x_p
\simeq K_N/K_Q
\simeq(S^1)^{\dim Q},
}
\]
where \(x_p\) is the unique point of the positive real part mapping to
\(p\).
\end{theorem}

\noindent
\textit{Reference.}
This is Cox--Little--Schenck, Proposition~12.2.3(b), combined with
Theorem~12.2.5 and Exercise~12.2.8.

\begin{proof}[Conceptual proof sketch]
The positive-part homeomorphism supplies the unique point \(x_p\) over
\(p\) with no angular coordinates. Polar decomposition of
\(O_Q\simeq(\CC^*)^{\dim Q}\) says that every other point over \(p\)
is obtained by restoring those angular coordinates, hence lies in
\(K_N\cdot x_p\). The stabilizer is \(K_Q\), so the fiber is
\(K_N/K_Q\simeq(S^1)^{\dim Q}\).
\end{proof}

\begin{corollary}[The polytope is the orbit space]
\label{cor:polytope-compact-orbit-space}
The moment map induces a homeomorphism
\[
\boxed{
\overline{\mu}_P:X_P/K_N\xrightarrow{\ \sim\ }P.
}
\]
\end{corollary}

\noindent
\textit{Reference.}
Cox--Little--Schenck, Proposition~12.2.3(a) identifies the compact-torus
quotient with the nonnegative part; Theorem~12.2.5 and
Exercise~12.2.8 identify that part with \(P\) via the moment map.

\begin{proof}[Conceptual proof sketch]
The preceding theorem says exactly that the fibers of \(\mu_P\) are
the \(K_N\)-orbits. Thus \(\mu_P\) induces a continuous bijection
\(X_P/K_N\to P\), which is a homeomorphism because the source is
compact and the target is Hausdorff.
\end{proof}

\begin{threedictionaries}
For a face \(Q\preceq P\) of codimension \(k\), let
\(p\in\relint(Q)\). The three versions of the
orbit--face dictionary are:

\medskip

\begin{center}
\small
\begin{tabularx}{\linewidth}{@{}
  >{\raggedright\arraybackslash}X
  >{\raggedright\arraybackslash}X
  >{\raggedright\arraybackslash}X@{}}
\toprule
\textbf{Polytope}
&
\textbf{Algebraic geometry}
&
\textbf{Compact-torus topology}
\\
\midrule
\(\operatorname{Int}(P)\)
&
dense orbit \(T_N\)
&
free orbit \(K_N\simeq(S^1)^n\)
\\[1mm]
facet \(F\)
&
dense orbit in the divisor \(D_F\)
&
one circle stabilizer; fiber \((S^1)^{n-1}\)
\\[1mm]
codimension-\(k\) face \(Q\)
&
codimension-\(k\) orbit closure \(X_Q\)
&
\(k\)-torus stabilizer; fiber \((S^1)^{n-k}\)
\\[1mm]
vertex \(v\)
&
torus-fixed point \(x_v\)
&
stabilizer \(K_N\); fiber is one point
\\
\bottomrule
\end{tabularx}
\end{center}
\end{threedictionaries}

\subsection{A toric surface over its polygon}

When \(\dim X_P=2\), the theorem can be seen at a glance. The compact
torus is
\[
K_N\simeq(S^1)^2.
\]
An interior point of the polygon has trivial stabilizer, so its fiber
is a two-dimensional torus. Along the relative interior of an edge,
one circle is a stabilizer and the fiber is the other circle. At a
vertex, both circle directions stabilize the point, so the fiber has
collapsed completely.

\begin{figure}[ht]
\centering
\begin{tikzpicture}[scale=0.9]
  % Interior panel.
  \begin{scope}[xshift=-4.7cm]
    \node[font=\bfseries] at (1.35,4.05) {interior};
    \draw[very thick,DictionaryBlue]
      (1.35,3.05) ellipse [x radius=0.67,y radius=0.25];
    \draw[very thick,DictionaryBlue]
      (1.35,3.05) ellipse [x radius=0.25,y radius=0.67];
    \node[right,DictionaryBlue] at (2.05,3.05) {\((S^1)^2\)};
    \draw[-{Latex[length=2.5mm]},thick,black!55]
      (1.35,2.35)--(1.35,1.65) node[midway,right] {\(\mu\)};
    \fill[DictionaryBlue!12]
      (0,0)--(2.7,0)--(2.3,1.45)--(0.45,1.75)--cycle;
    \draw[polytope,DictionaryBlue]
      (0,0)--(2.7,0)--(2.3,1.45)--(0.45,1.75)--cycle;
    \fill[DictionaryBlue] (1.35,0.82) circle (2.7pt);
  \end{scope}

  % Edge panel.
  \begin{scope}[xshift=0cm]
    \node[font=\bfseries] at (1.35,4.05) {edge};
    \draw[very thick,ExampleGold]
      (1.35,3.05) circle (0.53);
    \node[right,ExampleGold] at (1.98,3.05) {\(S^1\)};
    \draw[-{Latex[length=2.5mm]},thick,black!55]
      (1.35,2.35)--(1.35,1.65) node[midway,right] {\(\mu\)};
    \fill[DictionaryBlue!8]
      (0,0)--(2.7,0)--(2.3,1.45)--(0.45,1.75)--cycle;
    \draw[polytope,DictionaryBlue]
      (0,0)--(2.7,0)--(2.3,1.45)--(0.45,1.75)--cycle;
    \fill[ExampleGold] (2.5,0.72) circle (2.7pt);
  \end{scope}

  % Vertex panel.
  \begin{scope}[xshift=4.7cm]
    \node[font=\bfseries] at (1.35,4.05) {vertex};
    \fill[PausePurple] (1.35,3.05) circle (3.2pt);
    \node[right,PausePurple] at (1.57,3.05) {point};
    \draw[-{Latex[length=2.5mm]},thick,black!55]
      (1.35,2.35)--(1.35,1.65) node[midway,right] {\(\mu\)};
    \fill[DictionaryBlue!8]
      (0,0)--(2.7,0)--(2.3,1.45)--(0.45,1.75)--cycle;
    \draw[polytope,DictionaryBlue]
      (0,0)--(2.7,0)--(2.3,1.45)--(0.45,1.75)--cycle;
    \fill[PausePurple] (0,0) circle (2.7pt);
  \end{scope}
\end{tikzpicture}
\caption{The moment fibers of a toric surface. Two circle directions
survive over the polygon interior, one survives over an edge, and none
survive over a vertex. Thus the boundary records exactly where compact
torus directions collapse.}
\label{fig:toric-surface-over-moment-polygon}
\end{figure}

\begin{example}[Which circles collapse on a rectangle?]
\label{ex:circle-collapse-product-lines}
For \(X=\PP^1\times\PP^1\), write
\[
K_N=(S^1)_x\times(S^1)_y
\]
and use the square moment polytope \([0,1]^2\). The product formula from
the rectangle example in Section~8.1 gives
\[
\mu^{-1}(s,t)
=
\mu_{\PP^1}^{-1}(s)\times\mu_{\PP^1}^{-1}(t).
\]
Consequently:
\[
\begin{array}{c|c|c}
\textbf{position of }(s,t)
&
\textbf{stabilizer}
&
\textbf{fiber}
\\
\hline
0<s<1,\ 0<t<1
&
\{1\}
&
(S^1)_x\times(S^1)_y
\\
s\in\{0,1\},\ 0<t<1
&
(S^1)_x
&
(S^1)_y
\\
0<s<1,\ t\in\{0,1\}
&
(S^1)_y
&
(S^1)_x
\\
s,t\in\{0,1\}
&
(S^1)_x\times(S^1)_y
&
\text{point}
\end{array}
\]
The two vertical sides record the collapse of the \(x\)-circle, and
the two horizontal sides record the collapse of the \(y\)-circle.
At a corner both collapses occur simultaneously.
\end{example}

\subsection{Reconstructing the space from collapsing tori}

Over the interior, the moment map is a genuine product. Indeed, the
positive real slice and polar decomposition give
\[
\mu_P^{-1}(\operatorname{Int}P)
\simeq
\operatorname{Int}(P)\times K_N.
\]
The boundary prevents this product description from extending
unchanged to all of \(P\): the subgroup \(K_Q\) must be collapsed over
\(\relint(Q)\).

This can be packaged as a global topological model. Let
\[
s:P\longrightarrow(X_P)_{\geq0}
\]
be the inverse of the positive-real-part homeomorphism. For
\(p\in\relint(Q)\), impose the relation
\[
(p,k)\sim(p,k')
\quad\Longleftrightarrow\quad
k^{-1}k'\in K_Q.
\]

\begin{proposition}[The collapsing-torus model]
\label{prop:collapsing-torus-model}
The map
\[
P\times K_N\longrightarrow X_P,
\qquad
(p,k)\longmapsto k\cdot s(p),
\]
induces a homeomorphism
\[
\boxed{
X_P\simeq(P\times K_N)/{\sim}.
}
\]
Over \(\relint(Q)\), the relation collapses exactly the
subtorus \(K_Q\).
\end{proposition}

\noindent
\textit{Reference.}
This is the topological model in Cox--Little--Schenck,
equation~(12.2.7); its construction is Exercise~12.2.9.

\begin{proof}[Conceptual proof sketch]
The section \(s\) chooses the positive representative of each orbit,
and polar decomposition restores an arbitrary angular part
\(k\in K_N\). Over \(p\in\relint(Q)\), two angular parts
give the same point precisely when they differ by the stabilizer
\(K_Q\). This is exactly the relation \(\sim\); compactness then turns
the resulting continuous bijection into a homeomorphism.
\end{proof}

This model makes the slogan ``the polytope is the orbit space'' more
informative. The polytope specifies where the orbits live, while its
face normals specify which circle subgroups collapse. The total space
is obtained by starting with a full compact torus above the interior
and performing the prescribed collapses along the boundary.

\begin{warningbox}
The homeomorphism
\[
X_P/K_N\simeq P
\]
does \emph{not} say that \(X_P\) is homeomorphic to \(P\). The left-hand
side has already collapsed every compact-torus orbit. Nor does the
underlying topological shape of \(P\) alone remember which circle
subgroup collapses along each facet; that information comes from the
lattice normals.
\end{warningbox}

\subsection{What changes in the singular case?}

Neither Theorem~\ref{thm:moment-fibers-compact-orbits} nor
Corollary~\ref{cor:polytope-compact-orbit-space} requires \(X_P\) to be smooth.
For a singular projective toric variety, the orbit space is still the
polytope, and the fiber over \(\relint(Q)\) is still the
compact torus \(K_N/K_Q\).

What can fail is the local topology of the total space near a collapsed
fiber. For example, let \(P\) be a lattice polygon and let \(v\) be a
vertex. Suppose the normal cone at \(v\) is generated by primitive rays
\[
\sigma_v=\Cone{u_1,u_2}.
\]
Then
\[
d=|\det(u_1,u_2)|
\]
measures the local lattice index. If \(d=1\), the affine chart at the
fixed point is smooth and isomorphic to \(\CC^2\). If \(d>1\), the
fixed point is a cyclic quotient singularity. In both cases the moment
fiber over \(v\) is one point, and a neighborhood of \(v\) in the
orbit space still looks like a corner of a polygon.

Thus the orbit-space theorem records the stratification correctly but
does not, by itself, guarantee smooth local models. Smoothness requires
the familiar lattice-basis condition on every maximal cone.

\begin{observation}[A triangle need not mean \(\PP^2\)]
The moment polytope of a polarized weighted projective plane is also a
triangle. Over its interior sit two-dimensional compact tori, over its
edges sit circles, and over its vertices sit fixed points, exactly as
for \(\PP^2\). A non-unimodular normal cone at one vertex nevertheless
produces a quotient singularity there. The integral geometry of the
triangle, not its shape as an unlabelled disk, detects the difference.
\end{observation}

\begin{prediction}
Choose a generic vector \(u\in N_{\RR}\) and use it to define the height
function
\[
h_u(m)=\langle m,u\rangle
\]
on \(P\). Orient every edge in the direction of increasing height.

At a vertex \(v\), the fiber has collapsed to a fixed point \(x_v\).
How might the one-parameter subgroup determined by \(u\) use the
oriented edges at \(v\) to select a whole region of \(X_P\) flowing
toward \(x_v\)? What should the dimension of that region be?
\end{prediction}

\begin{guidingquestion}
Can a generic height function decompose a smooth projective toric
variety into affine cells indexed by the vertices of \(P\), so that
the Betti numbers can be read directly from the oriented polytope?
\end{guidingquestion}

\section{One-parameter flows, cells, and Betti numbers}
\label{sec:one-parameter-flows-cells-betti}
\index{one-parameter flow}
\index{attracting cell}
\index{cell decomposition!toric variety}
\index{Betti number}
\index{h-vector}

The orbit-space picture tells us where the compact-torus directions
collapse. We now let the noncompact part of the algebraic torus move
points toward those collapsed fibers. A generic one-parameter subgroup
will divide the whole variety into attracting regions, one for each
torus fixed point. When the variety is smooth, these regions are affine
spaces. Thus a height function on the polytope becomes a tool for
computing topology.

Throughout this section, \(P\subset M_{\RR}\) is a full-dimensional
lattice polytope and \(X_P\) is smooth and projective of dimension \(n\).

\subsection{One vector, two kinds of motion}

Choose a lattice vector \(u\in N\). It defines a one-parameter subgroup
\[
\lambda_u:\CC^*\longrightarrow T_N,
\qquad
t\longmapsto\lambda_u(t),
\]
characterized by
\[
\chi^m(\lambda_u(t))=t^{\langle m,u\rangle}
\qquad(m\in M).
\]
The same vector defines a linear height function
\[
h_u:P\longrightarrow\RR,
\qquad
h_u(m)=\langle m,u\rangle.
\]

We call \(u\) \emph{generic for \(P\)} if it is not perpendicular to
any edge of \(P\). Equivalently, \(h_u\) takes distinct values at the
two endpoints of every edge. For such a vector, orient each edge from
lower height to higher height.

\begin{observation}[The flow is visible on the polytope]
On the dense torus, write a point as \(k\exp(-w)\), with
\(k\in K_N\) and \(w\in N_{\RR}\). Multiplication by
\(\lambda_u(t)\), for positive real \(t\), translates the real
coordinate in the direction \(-\log(t)u\). Under the moment map, the
resulting path moves monotonically in the direction of decreasing
\(h_u\) as \(t\to0\). Its limiting vertex is the minimum of \(h_u\)
on the appropriate face.
\end{observation}

The convention in this observation is worth remembering: the limit
\(t\to0\) corresponds to movement toward \emph{smaller} height. Replacing
\(u\) by \(-u\), or taking \(t\to\infty\), reverses every arrow and
interchanges attracting and repelling sets.

\subsection{Attracting sets at the fixed points}

The vertices of \(P\) correspond to the torus fixed points of \(X_P\).
Write \(x_v\) for the fixed point corresponding to a vertex \(v\). Its
attracting set is
\[
C_v
=
\left\{
x\in X_P:
\lim_{t\to0}\lambda_u(t)\cdot x=x_v
\right\}.
\]

To read its dimension, look only at the oriented edges incident to
\(v\). Let
\[
d_u(v)
=
\#\{\text{edges directed away from }v\}.
\]
Thus \(d_u(v)\) is the number of edge directions along which height
increases from \(v\).

\begin{theorem}[Attracting cells from oriented edges]
\label{thm:toric-attracting-cells}
Let \(X_P\) be a smooth projective toric variety and let \(u\in N\) be
generic for \(P\). Then
\[
X_P=\coprod_{v\in\operatorname{Vert}(P)}C_v,
\qquad
C_v\simeq\AAA^{d_u(v)}_{\CC}.
\]
Moreover, the cells can be ordered so that successive unions are
closed. Consequently they give a CW decomposition with one real
\(2d_u(v)\)-dimensional cell for each vertex \(v\).
\end{theorem}

\noindent
\textit{Reference.}
This is the toric form of the Bia\l ynicki--Birula decomposition.
Cox--Little--Schenck do not formulate the decomposition in this
language; their orbit filtration in \S12.3 gives the corresponding
Betti-number count in Theorem~12.3.12. The local weight calculation
uses the smooth affine charts described in \S3.3.

\begin{proof}[Conceptual proof sketch]
At a fixed point \(x_v\), smoothness identifies the invariant affine
chart with \(\AAA^n\), one coordinate for each edge leaving \(v\). If
the corresponding primitive edge direction is \(e\in M\), then
\(\lambda_u(t)\) acts on that coordinate with weight
\(\langle e,u\rangle\). A coordinate may be arbitrary when this weight
is positive, since multiplication by \(t^{\langle e,u\rangle}\) sends
it to zero as \(t\to0\); it must vanish when the weight is negative.
Hence the attracting set is an affine coordinate subspace of dimension
\(d_u(v)\). Projectivity orders these strata by height and makes their
successive unions closed, producing the asserted cell decomposition.
\end{proof}

\begin{warningbox}
The word \emph{generic} is essential. If \(u\) is perpendicular to an
edge, the fixed locus of \(\lambda_u\) can contain a positive-dimensional
toric subvariety rather than isolated fixed points. The resulting
Bia\l ynicki--Birula strata are then affine bundles over fixed components,
not necessarily affine spaces indexed by vertices.
\end{warningbox}

\subsection{Reading the cells from a polygon}

For a toric surface, each vertex has two incident edges. After orienting
them, its attracting cell has dimension \(0\), \(1\), or \(2\), according
to whether zero, one, or two arrows leave the vertex.

\begin{figure}[ht]
\centering
\begin{tikzpicture}[scale=1.0,
  flowarrow/.style={-{Latex[length=2.5mm]},very thick,DictionaryBlue}]
  \coordinate (a) at (0,0);
  \coordinate (b) at (4.4,0.25);
  \coordinate (c) at (5.15,2.65);
  \coordinate (d) at (2.15,3.35);
  \coordinate (e) at (-0.45,1.9);

  \fill[DictionaryBlue!7] (a)--(b)--(c)--(d)--(e)--cycle;
  \draw[polytope,black!60] (a)--(b)--(c)--(d)--(e)--cycle;

  \draw[flowarrow] (a)--(b);
  \draw[flowarrow] (b)--(c);
  \draw[flowarrow] (e)--(d);
  \draw[flowarrow] (d)--(c);
  \draw[flowarrow] (a)--(e);

  \fill[PausePurple] (a) circle (2.7pt)
    node[below left] {\(\dim C_v=2\)};
  \fill[PausePurple] (b) circle (2.7pt)
    node[below right] {\(1\)};
  \fill[PausePurple] (c) circle (2.7pt)
    node[right] {\(0\)};
  \fill[PausePurple] (d) circle (2.7pt)
    node[above] {\(1\)};
  \fill[PausePurple] (e) circle (2.7pt)
    node[left] {\(1\)};

  \draw[-{Latex[length=3mm]},thick,ExampleGold]
    (0.55,-0.8)--(3.35,-0.8)
    node[midway,below] {increasing \(h_u\)};
\end{tikzpicture}
\caption{A generic height orients the edges. The label at a vertex is
the number of arrows leaving it, hence the complex dimension of its
attracting cell. This pentagon gives one \(0\)-cell, three
\(1\)-cells, and one \(2\)-cell.}
\label{fig:oriented-polygon-attracting-cells}
\end{figure}

\begin{example}[The projective plane]
For the standard triangle, a generic height orders the three vertices.
At the minimum, both edges point outward; at the middle vertex, one
points in and one points out; at the maximum, both point inward. Thus
\[
\PP^2
=
\AAA^2\sqcup\AAA^1\sqcup\AAA^0.
\]
The three cells have real dimensions \(4,2,0\), exactly as in the
usual CW decomposition of projective space.
\end{example}

\begin{example}[The product of two projective lines]
Choose a generic diagonal height on the square. The numbers of outgoing
edges at its four vertices are
\[
2,\quad1,\quad1,\quad0.
\]
Therefore
\[
\PP^1\times\PP^1
=
\AAA^2\sqcup\AAA^1\sqcup\AAA^1\sqcup\AAA^0.
\]
The two one-dimensional cells already foreshadow the two independent
degree-two classes coming from the two projections.
\end{example}

\begin{example}[Cutting off a corner]
The moment polygon of \(\Bl_p(\PP^2)\) is obtained by cutting one corner
off a triangle. A generic orientation of the resulting quadrilateral
has outgoing-edge counts
\[
2,\quad1,\quad1,\quad0.
\]
Thus the blowup has the same number of cells in each dimension as
\(\PP^1\times\PP^1\), even though the two toric surfaces have different
intersection forms. Additive topology sees the extra degree-two class;
the ring structure will distinguish how that class intersects itself.
\end{example}

\subsection{Odd cohomology vanishes}

Because every cell has even real dimension, the cellular cochain
complex has no groups in odd degree. There are also no differentials
between even-dimensional cells: a cellular differential changes degree
by one.

\begin{corollary}[Betti numbers from the oriented polytope]
\label{cor:betti-oriented-polytope}
For a smooth projective toric variety,
\[
H^{2k+1}(X_P;\ZZ)=0
\qquad\text{and}\qquad
b_{2k}(X_P)=\#\{v\in\operatorname{Vert}(P):d_u(v)=k\}.
\]
In particular, the integral cohomology groups are free abelian and
\[
\chi(X_P)
=|\operatorname{Vert}(P)|
=|\Sigma(n)|.
\]
\end{corollary}

\noindent
\textit{Reference.}
The vanishing of odd cohomology and the combinatorial formula for the
even Betti numbers are Cox--Little--Schenck, Theorems~12.3.11
and~12.3.12. The Euler-characteristic formula is Theorem~12.3.9.

\begin{proof}[Conceptual proof sketch]
Theorem~\ref{thm:toric-attracting-cells} supplies one even-dimensional
cell for each vertex. Cellular cohomology therefore has no odd groups
and no nonzero differentials. A cell \(C_v\simeq\CC^k\) contributes one
generator in degree \(2k\). Summing the ranks gives one contribution per
vertex, so the Euler characteristic is the number of vertices. Under
the orbit--cone dictionary, vertices of \(P\) correspond to maximal
cones of \(\Sigma\).
\end{proof}

\begin{observation}[Why the answer cannot depend on \(u\)]
Changing the generic height may reverse some arrows and move cells from
one fixed point to another, but Corollary~\ref{cor:betti-oriented-polytope}
shows that the number of vertices with \(k\) outgoing edges is invariant.
The individual cells depend on \(u\); their number in each dimension is
a topological invariant of \(X_P\).
\end{observation}

\subsection{The \texorpdfstring{\(h\)}{h}-vector packages the count}

The preceding examples suggest attaching a polynomial to a simple
polytope by counting vertices according to the number of outgoing
edges. This is one of the standard descriptions of the \(h\)-vector.

\begin{definition}[The \(h\)-vector of a simple polytope]
Let \(P\) be a simple \(n\)-polytope and choose a generic height
function. Define
\[
h_k(P)=\#\{v\in\operatorname{Vert}(P):d_u(v)=k\},
\qquad 0\leq k\leq n.
\]
The vector
\[
h(P)=(h_0,h_1,\ldots,h_n)
\]
is independent of the chosen generic height and is called the
\emph{\(h\)-vector} of \(P\).
\end{definition}

For readers who prefer a definition using only face numbers, let
\(f_i\) denote the number of \(i\)-dimensional faces of \(P\), with
\(f_{-1}=1\). The same numbers are characterized by
\[
\sum_{k=0}^{n}h_k t^{n-k}
=
\sum_{i=0}^{n}f_{i-1}(t-1)^{n-i}.
\]
The oriented-vertex description explains geometrically why these
apparently combinatorial coefficients must be nonnegative.

\begin{theorem}[The \(h\)-vector is the Betti vector]
\label{thm:h-vector-betti-vector}
If \(P\) is a smooth lattice polytope, then
\[
\boxed{
b_{2k}(X_P)=h_k(P),
\qquad
b_{2k+1}(X_P)=0.
}
\]
Equivalently, the Poincar\'e polynomial is
\[
\operatorname{Poin}_{X_P}(q)
=
\sum_{k=0}^{n}h_k(P)q^{2k}.
\]
\end{theorem}

\noindent
\textit{Reference.}
This is the smooth projective case of Cox--Little--Schenck,
Theorem~12.3.12.

\begin{proof}[Conceptual proof sketch]
By definition, \(h_k(P)\) counts precisely the vertices whose attracting
cells have complex dimension \(k\). The result is therefore a
repackaging of Corollary~\ref{cor:betti-oriented-polytope}.
\end{proof}

Poincar\'e duality now translates into the symmetry
\[
h_k(P)=h_{n-k}(P).
\]
Thus the topology of \(X_P\) gives a geometric explanation for the
Dehn--Sommerville symmetry of the \(h\)-vector in the smooth case.

\begin{example}[Three familiar \(h\)-vectors]
The cell counts above give
\[
\begin{array}{c|c|c}
P & h(P) & \operatorname{Poin}_{X_P}(q)\\
\hline
\text{triangle} & (1,1,1) & 1+q^2+q^4\\
\text{square} & (1,2,1) & 1+2q^2+q^4\\
\text{corner-cut triangle} & (1,2,1) & 1+2q^2+q^4
\end{array}
\]
The last two rows again show that Betti numbers are coarser than the
toric or algebraic geometry.
\end{example}

\subsection{Toric surfaces and blowups}

Every convex \(r\)-gon has one vertex with no outgoing edge, one vertex
with two outgoing edges, and \(r-2\) vertices with one outgoing edge.
Indeed, the first two are the maximum and minimum of a generic height;
every other vertex has one neighbor above it and one below it.
Therefore its \(h\)-vector is
\[
(1,r-2,1).
\]

\begin{corollary}[Betti numbers of a smooth complete toric surface]
\label{cor:betti-smooth-toric-surface}
Let \(X_\Sigma\) be a smooth complete toric surface and let
\(r=|\Sigma(1)|\). Then
\[
b_0=b_4=1,
\qquad
b_2=r-2,
\qquad
b_1=b_3=0.
\]
Each toric blowup at a fixed point adds one ray to the fan and therefore
increases \(b_2\) by one.
\end{corollary}

\begin{proof}[Conceptual proof sketch]
A complete two-dimensional fan has as many rays as maximal cones, so
its moment polygon has \(r\) vertices. The oriented polygon count gives
one cell in complex dimensions \(0\) and \(2\), and \(r-2\) cells in
complex dimension \(1\). A toric blowup star-subdivides one maximal cone,
adding one ray and replacing one polygon vertex by an edge.
\end{proof}

\begin{example}[The first blowup of \(\PP^2\)]
The fan of \(\PP^2\) has three rays, so \(b_2=1\). Blowing up a torus
fixed point adds the exceptional ray and gives
\[
b_2(\Bl_p\PP^2)=4-2=2.
\]
One degree-two class is inherited from a line in \(\PP^2\), while the
other is represented by the exceptional divisor. The cell decomposition
detects that there are two classes, but it does not yet compute
\(H^2=1\), \(E^2=-1\), or \(H\cdot E=0\). For those products we need the
cohomology ring.
\end{example}

\begin{pauseandthink}
Choose a generic height on each polygon below and verify the stated
Betti numbers directly from the oriented edges.
\begin{enumerate}
\item A triangle gives \((b_0,b_2,b_4)=(1,1,1)\).
\item A quadrilateral gives \((1,2,1)\).
\item A hexagon gives \((1,4,1)\).
\end{enumerate}
Why can no choice of generic height change these counts?
\end{pauseandthink}

\begin{threedictionaries}
\begin{center}
\begin{tabularx}{\textwidth}{>{\bfseries}lXXX}
\toprule
& Polytope & One-parameter flow & Topology \\
\midrule
Indexing object
& vertex \(v\)
& fixed point \(x_v\)
& cell \(C_v\)
\\
Local number
& outgoing edges \(d_u(v)\)
& positive weights at \(x_v\)
& \(\dim_{\CC}C_v\)
\\
Global count
& \(h_k(P)\)
& fixed points with \(k\) positive weights
& \(b_{2k}(X_P)\)
\\
Parity
& ---
& affine attracting spaces
& \(H^{\mathrm{odd}}=0\)
\\
\bottomrule
\end{tabularx}
\end{center}
\end{threedictionaries}

\begin{prediction}
The cell decomposition determines the additive groups
\(H^*(X_P;\ZZ)\), but not their multiplication. The codimension-one
cells suggest that degree-two cohomology should be generated by
invariant divisors. Which two families of relations among those divisors
should come from empty intersections and principal divisors?
\end{prediction}

\begin{guidingquestion}
Can the fan recover not only the Betti numbers but the entire integral
cohomology ring of a smooth complete toric variety?
\end{guidingquestion}

\section{The cohomology ring from the fan}
\label{sec:cohomology-ring-from-fan}
\index{cohomology ring!toric presentation}
\index{Stanley--Reisner ideal}
\index{Jurkiewicz--Danilov presentation}
\index{intersection number!from cohomology}

The cell decomposition from the previous section determines the
additive groups
\[
H^k(X_\Sigma;\ZZ).
\]
It does not yet describe the cup product. Two smooth toric surfaces can
have the same number of cells in every dimension while their
degree-two classes multiply differently.

The fan contains exactly the extra information we need. Its rays index
the invariant divisors, its cones record which of those divisors meet,
and its lattice vectors record the linear equivalences among them.
Those two kinds of data become the two kinds of relations in the
cohomology ring.

\subsection{Invariant divisors as degree-two classes}

Let \(X_\Sigma\) be a smooth complete toric variety. For every ray
\(\rho\in\Sigma(1)\), let
\[
D_\rho=V(\rho)
\]
be the corresponding invariant prime divisor. Its cohomology class is
\[
x_\rho
:=
c_1\bigl(\mathcal O_{X_\Sigma}(D_\rho)\bigr)
=
\operatorname{PD}[D_\rho]
\in H^2(X_\Sigma;\ZZ).
\]
Here \(\operatorname{PD}\) denotes Poincar\'e duality. Thus the symbol
\([D_\rho]\), when used in cohomology, means the degree-two class
Poincar\'e dual to the divisor rather than merely its class in
\(\Pic(X_\Sigma)\).

Cup product turns intersections into multiplication. If distinct rays
\(\rho_1,\ldots,\rho_k\) span a cone \(\sigma\), smoothness implies
that the corresponding boundary divisors meet transversely and
\[
x_{\rho_1}\cdots x_{\rho_k}
=
\operatorname{PD}[V(\sigma)].
\]
If the rays do not span a cone, the common intersection is empty and
the product vanishes. This gives the first family of relations.

\begin{definition}[The two fan ideals]
\label{def:two-fan-ideals}
Let
\[
S^{\mathrm{top}}_\Sigma
=
\ZZ[x_\rho:\rho\in\Sigma(1)],
\qquad
\deg x_\rho=2.
\]

The \emph{Stanley--Reisner ideal} is
\[
I_\Sigma
=
\left\langle
x_{\rho_1}\cdots x_{\rho_k}:
\rho_1,\ldots,\rho_k
\text{ do not span a cone of }\Sigma
\right\rangle,
\]
where the rays in each generator are distinct.

For \(m\in M\), set
\[
\theta_m
=
\sum_{\rho\in\Sigma(1)}
\langle m,u_\rho\rangle x_\rho.
\]
The \emph{linear-relation ideal} is
\[
J_\Sigma
=
\langle\theta_m:m\in M\rangle.
\]
It is enough to use the elements \(\theta_{m_1},\ldots,
\theta_{m_n}\) coming from a basis of \(M\).
\end{definition}

The second ideal is familiar from Chapter~4. Every character has
principal divisor
\[
\operatorname{div}(\character{m})
=
\sum_{\rho\in\Sigma(1)}
\langle m,u_\rho\rangle D_\rho.
\]
A principal divisor defines the trivial line bundle, so its first Chern
class vanishes. Therefore
\[
\theta_m
=
\sum_\rho\langle m,u_\rho\rangle x_\rho
=0
\qquad\text{in }H^2(X_\Sigma;\ZZ).
\]

We have now found an evident graded ring homomorphism
\[
\Phi:
\frac{\ZZ[x_\rho:\rho\in\Sigma(1)]}
     {I_\Sigma+J_\Sigma}
\longrightarrow
H^*(X_\Sigma;\ZZ),
\qquad
x_\rho\longmapsto\operatorname{PD}[D_\rho].
\]
The remarkable statement is that there are no additional generators
or relations.

\begin{theorem}[Jurkiewicz--Danilov presentation]
\label{thm:jurkiewicz-danilov-presentation}
Let \(X_\Sigma\) be a smooth complete toric variety. Then
\[
\boxed{
H^*(X_\Sigma;\ZZ)
\simeq
\frac{\ZZ[x_\rho:\rho\in\Sigma(1)]}
     {I_\Sigma+J_\Sigma},
\qquad
\deg x_\rho=2.
}
\]
Under this isomorphism, \(x_\rho\) is the cohomology class of the
invariant divisor \(D_\rho\).
\end{theorem}

\noindent
\textit{Reference.}
This is Cox--Little--Schenck, Theorem~12.4.4, where it is called the
Jurkiewicz--Danilov theorem. The rational presentation for complete
simplicial toric varieties is their Theorem~12.4.1.

\begin{proof}[Conceptual proof sketch]
The divisor classes give a map from the polynomial ring to cohomology.
Boundary incidence puts \(I_\Sigma\) in its kernel, and principal
divisors put \(J_\Sigma\) in its kernel.

The conceptual reason these are all the relations comes from
equivariant cohomology. Before forgetting the compact-torus action, the
equivariant cohomology ring is the Stanley--Reisner ring
\[
H^*_{K_N}(X_\Sigma;\ZZ)
\simeq
\frac{\ZZ[x_\rho]}{I_\Sigma}.
\]
The coefficient ring \(H^*(BK_N;\ZZ)\simeq\operatorname{Sym}_{\ZZ}(M)\)
acts, with our principal-divisor convention, by sending
\(m\in M\) to the linear form \(-\theta_m\). Passing
from equivariant to ordinary cohomology sets all positive-degree
classes from \(BK_N\) equal to zero, which imposes exactly the ideal
\(J_\Sigma\).

Smoothness and completeness ensure that this passage introduces no
hidden integral torsion or extra relations. In the projective case,
the even-dimensional cell decomposition from Section~8.3 is the
geometric shadow of this freeness.
\end{proof}

\begin{warningbox}
Both hypotheses matter. For a complete simplicial fan, the analogous
presentation holds with rational coefficients, but integral
cohomology can be subtler. For a nonsimplicial singular toric variety,
ordinary cohomology need not be generated by divisor classes at all.
The integral formula in
\cref{thm:jurkiewicz-danilov-presentation} is the smooth complete
statement.
\end{warningbox}

\subsection{Cox variables and cohomology variables}

The polynomial ring in
\cref{thm:jurkiewicz-danilov-presentation} looks deliberately similar
to the Cox ring from Chapter~6. The variables have the same indexing
set, but they play different roles.

\begin{dictionarytable}
\begin{center}
\small
\begin{tabularx}{\linewidth}{@{}
  >{\bfseries\raggedright\arraybackslash}p{0.21\linewidth}
  >{\raggedright\arraybackslash}X
  >{\raggedright\arraybackslash}X@{}}
\toprule
& \textbf{Cox construction} & \textbf{Cohomology ring} \\
\midrule
Coefficient ring
& \(\CC\)
& \(\ZZ\)
\\
Variable
& coordinate \(z_\rho\) cutting out \(D_\rho\)
& class \(x_\rho=\operatorname{PD}[D_\rho]\)
\\
Degree
& \(\deg z_\rho=[D_\rho]\in\Cl(X_\Sigma)\)
& \(\deg x_\rho=2\)
\\
Fan nonfaces
& the irrelevant ideal records excluded coordinate strata
& \(I_\Sigma\) kills empty divisor intersections
\\
Principal divisors
& determine the class-group grading
& generate the linear ideal \(J_\Sigma\)
\\
Final operation
& remove \(V(B(\Sigma))\), then take a quotient
& quotient the polynomial ring by \(I_\Sigma+J_\Sigma\)
\\
\bottomrule
\end{tabularx}
\end{center}
\end{dictionarytable}

Thus the two rings are not the same ring in different notation. The
Cox ring uses variables as coordinates before taking a geometric
quotient. The cohomology ring uses variables as divisor classes after
the variety has been constructed.

\subsection{A four-step computation}

For a smooth complete fan, the presentation can be calculated by a
short procedure.

\begin{enumerate}[label=\textbf{Step \arabic*.},wide=0pt]
\item List the rays \(\rho\in\Sigma(1)\) and introduce one variable
\(x_\rho\) for each.

\item Find the minimal collections of rays that do not span cones.
Their squarefree monomials generate \(I_\Sigma\).

\item Choose a basis of \(M\). Pair its elements with the primitive
ray generators to obtain generators of \(J_\Sigma\).

\item Use the linear relations to eliminate variables. The remaining
monomial relations reveal the cup products.
\end{enumerate}

We now carry out all four steps for the three running examples.

\subsection{The projective plane}

\begin{example}[The cohomology ring of \(\PP^2\)]
\label{ex:cohomology-ring-projective-plane}
Use the primitive ray generators
\[
u_0=-e_1-e_2,
\qquad
u_1=e_1,
\qquad
u_2=e_2.
\]
Every pair spans a cone, but the three rays together do not. Hence
\[
I_\Sigma=\langle x_0x_1x_2\rangle.
\]

For the dual basis \(m_1,m_2\) of \(M\), the two principal-divisor
relations are
\[
-x_0+x_1=0,
\qquad
-x_0+x_2=0.
\]
Therefore
\[
J_\Sigma
=
\langle x_1-x_0,\ x_2-x_0\rangle.
\]
If \(H\) denotes the common class of the three coordinate lines, then
\[
x_0=x_1=x_2=H.
\]
The Stanley--Reisner relation becomes \(H^3=0\), and the theorem gives
\[
\boxed{
H^*(\PP^2;\ZZ)
\simeq
\frac{\ZZ[H]}{\langle H^3\rangle},
\qquad
\deg H=2.
}
\]
As an abelian group, the ring has basis
\[
1,\quad H,\quad H^2,
\]
in degrees \(0,2,4\). The product \(H^2\) is the class of the
intersection point of two distinct coordinate lines.
\end{example}

\subsection{The product of two projective lines}

\begin{example}[The cohomology ring of \(\PP^1\times\PP^1\)]
\label{ex:cohomology-ring-product-lines}
Order the rays cyclically as
\[
u_1=e_1,
\qquad
u_2=e_2,
\qquad
u_3=-e_1,
\qquad
u_4=-e_2.
\]
The opposite pairs do not span cones, while every adjacent pair does.
Thus
\[
I_\Sigma
=
\langle x_1x_3,\ x_2x_4\rangle.
\]
The character relations are
\[
x_1-x_3=0,
\qquad
x_2-x_4=0,
\]
so
\[
J_\Sigma
=
\langle x_1-x_3,\ x_2-x_4\rangle.
\]

Set
\[
A=x_1=x_3,
\qquad
B=x_2=x_4.
\]
Geometrically, \(A\) and \(B\) are the classes of fibers of the two
projections. The two nonface relations become
\[
A^2=0,
\qquad
B^2=0.
\]
Consequently,
\[
\boxed{
H^*(\PP^1\times\PP^1;\ZZ)
\simeq
\frac{\ZZ[A,B]}{\langle A^2,B^2\rangle},
\qquad
\deg A=\deg B=2.
}
\]
An additive basis is
\[
1,\quad A,\quad B,\quad AB.
\]
The top class \(AB\) is represented by the fixed point where one
fiber from each ruling meets. The equations \(A^2=B^2=0\) say that two
distinct fibers from the same ruling are disjoint.
\end{example}

\subsection{The blowup of the projective plane}

\begin{example}[The cohomology ring of \(\Bl_p(\PP^2)\)]
\label{ex:cohomology-ring-blowup-projective-plane}
Blow up the fixed point corresponding to \(\Cone{e_1,e_2}\). As in
Chapters~4 and~5, order the rays cyclically as
\[
u_0=-e_1-e_2,
\qquad
u_1=e_1,
\qquad
u_E=e_1+e_2,
\qquad
u_2=e_2.
\]
The two nonadjacent pairs are
\[
\{\rho_0,\rho_E\}
\qquad\text{and}\qquad
\{\rho_1,\rho_2\}.
\]
Therefore
\[
I_\Sigma
=
\langle x_0x_E,\ x_1x_2\rangle.
\]

Pairing the ray generators with the dual basis of \(M\) gives
\[
-x_0+x_1+x_E=0,
\qquad
-x_0+x_E+x_2=0.
\]
Hence
\[
J_\Sigma
=
\langle-x_0+x_1+x_E,\ -x_0+x_E+x_2\rangle.
\]
Before eliminating variables, the theorem reads
\[
H^*(\Bl_p(\PP^2);\ZZ)
\simeq
\frac{\ZZ[x_0,x_1,x_E,x_2]}
{\langle
x_0x_E,\ x_1x_2,\
-x_0+x_1+x_E,\ -x_0+x_E+x_2
\rangle}.
\]

Now put
\[
H=x_0,
\qquad
E=x_E.
\]
The linear relations say
\[
x_1=x_2=H-E.
\]
The two Stanley--Reisner relations become
\[
HE=0,
\qquad
(H-E)^2=0.
\]
Since \(HE=0\), the second equation is equivalent to
\[
H^2+E^2=0.
\]
We obtain the two useful forms of the answer
\[
\boxed{
H^*(\Bl_p(\PP^2);\ZZ)
\simeq
\frac{\ZZ[H,E]}{\langle HE,(H-E)^2\rangle}
\simeq
\frac{\ZZ[H,E]}{\langle HE,H^2+E^2\rangle},
\quad
\deg H=\deg E=2.
}
\]

The classes have their familiar meanings. The divisor \(D_0\) misses
the blown-up point, so \(H\) is the pullback of the line class. The new
divisor is the exceptional curve \(E\), while
\[
x_1=x_2=H-E
\]
are the classes of the strict transforms of the two coordinate lines
through the blown-up point.

An additive basis is
\[
1,\quad H,\quad E,\quad H^2.
\]
Indeed, \(E^2=-H^2\), and every product of three positive-degree
generators vanishes. The class \(H^2=x_0x_1\) is represented by a
torus fixed point.
\end{example}

\begin{surfacelaboratory}
The square and the corner-cut triangle gave the same Betti numbers:
\[
b_0=b_4=1,
\qquad
b_2=2.
\]
Their cohomology rings are nevertheless different.

On \(\PP^1\times\PP^1\), every degree-two class has the form
\(aA+bB\), and
\[
(aA+bB)^2=2ab\,AB.
\]
Thus every square is an even multiple of the top-degree generator. On
\(\Bl_p(\PP^2)\), however,
\[
E^2=-H^2,
\]
an odd multiple of a top-degree generator. Hence the two integral
cohomology rings cannot be isomorphic as graded rings.

Additive topology counts the two degree-two classes. Multiplication
remembers the parity of their intersection form.
\end{surfacelaboratory}

\subsection{From ring products to intersection numbers}

If \(X_\Sigma\) has complex dimension \(n\), then
\[
H^{2n}(X_\Sigma;\ZZ)\simeq\ZZ.
\]
After multiplying \(n\) divisor classes, evaluation on the fundamental
class produces an integer:
\[
\int_{X_\Sigma}
x_{\rho_1}\cdots x_{\rho_n}
\in\ZZ.
\]
This is the intersection number
\[
D_{\rho_1}\cdots D_{\rho_n}.
\]
For a smooth maximal cone
\(\sigma=\Cone{u_{\rho_1},\ldots,u_{\rho_n}}\), the corresponding
divisors meet transversely at one fixed point, so
\[
\int_{X_\Sigma}x_{\rho_1}\cdots x_{\rho_n}=1.
\]
In the surface examples, this normalizes
\[
\int_{\PP^2}H^2=1,
\qquad
\int_{\PP^1\times\PP^1}AB=1,
\qquad
\int_{\Bl_p(\PP^2)}H^2=1.
\]

\begin{threedictionaries}
\begin{center}
\small
\begin{tabularx}{\linewidth}{@{}
  >{\raggedright\arraybackslash}X
  >{\raggedright\arraybackslash}X
  >{\raggedright\arraybackslash}X@{}}
\toprule
\textbf{Fan}
& \textbf{Divisor geometry}
& \textbf{Cohomology ring}
\\
\midrule
rays not contained in a cone
& empty common intersection
& monomial in \(I_\Sigma\)
\\
character pairing with ray generators
& principal divisor
& linear form in \(J_\Sigma\)
\\
smooth maximal cone
& transverse fixed point
& top product equal to the point class
\\
\bottomrule
\end{tabularx}
\end{center}
\end{threedictionaries}

\begin{prediction}
Let the rays of a smooth complete toric surface be ordered cyclically,
and let \(x_i=[D_i]\). The Stanley--Reisner relations immediately give
\[
x_ix_j=0
\]
for nonadjacent rays, while adjacent rays give the point class. The
only products still to determine are the self-intersections
\[
x_i^2.
\]
How should the lattice relation among
\(u_{i-1},u_i,u_{i+1}\) determine this square?
\end{prediction}

\clearpage
\section*{What this chapter has accomplished}
\addcontentsline{toc}{section}{What this chapter has accomplished}
\markboth{What this chapter has accomplished}{What this chapter has accomplished}

Over \(\CC\), the algebraic torus contains its compact part
\[
K_N\simeq(S^1)^n\subseteq T_N\simeq(\CC^*)^n.
\]
For a projective toric embedding with lattice weights
\(A=\{m_0,\ldots,m_r\}\), normalized squared coordinate magnitudes
produce the moment map
\[
\mu_A([z_0:\cdots:z_r])
=
\frac{\sum_i|z_i|^2m_i}{\sum_i|z_i|^2}.
\]
Its coefficients form a convex combination, and the positive real
part of the toric variety shows that every point of the convex hull is
attained:
\[
\boxed{\mu_A(X_A)=\Conv(A).}
\]

The moment map does more than record an image. Its fibers are exactly
the compact-torus orbits. If \(p\in\relint(Q)\) for a
face \(Q\preceq P\), then
\[
\mu_P^{-1}(p)
\simeq K_N/K_Q
\simeq(S^1)^{\dim Q},
\]
where the stabilizer \(K_Q\) has dimension \(\codim Q\). Consequently,
\[
\boxed{X_P/K_N\simeq P.}
\]
The polytope is therefore a stratified orbit space: full compact tori
sit over its interior, circle directions collapse along facets, and
the whole compact torus collapses over a vertex.

A generic vector \(u\in N_\RR\) turns this orbit picture into a cell
decomposition. The height function
\[
h_u(m)=\langle m,u\rangle
\]
orients every edge of \(P\), and the fixed point over a vertex \(v\)
has an attracting cell
\[
C_v\simeq\CC^{d_u(v)},
\]
where \(d_u(v)\) is the number of edges leaving \(v\). All cells have
even real dimension, so
\[
H^{2k+1}(X_P;\ZZ)=0,
\qquad
b_{2k}(X_P)=h_k(P).
\]
In particular,
\[
\chi(X_P)=|\operatorname{Vert}(P)|=|\Sigma(n)|.
\]
For a smooth complete toric surface with \(s\) rays this specializes to
\[
b_0=b_4=1,
\qquad
b_2=s-2,
\qquad
b_1=b_3=0.
\]

Finally, the fan recovers not only the additive groups but their
multiplication. The invariant divisors generate integral cohomology,
and the two familiar families of toric relations give the complete
presentation
\[
\boxed{
H^*(X_\Sigma;\ZZ)
\simeq
\frac{\ZZ[x_\rho:\rho\in\Sigma(1)]}
     {I_\Sigma+J_\Sigma},
\qquad
\deg x_\rho=2.
}
\]
The Stanley--Reisner ideal \(I_\Sigma\) records empty intersections,
while \(J_\Sigma\) records principal divisors. Thus the same fan data
now govern three layers of topology:
\[
\boxed{
\begin{gathered}
\text{faces and compact-torus fibers},
\qquad
\text{cells and Betti numbers},
\\[-1mm]
\text{divisor intersections and the cohomology ring}.
\end{gathered}}
\]
The challenge exercise completes the surface story by extracting the
entire intersection form from a cyclic fan and asking how much of the
Hirzebruch parameter survives after the toric structure is forgotten.

\begin{exercises}[Comprehension exercises]
\markboth{Comprehension exercises}{Comprehension exercises}

\item For \(\PP^1\), use the weights \(0,1\in\ZZ\) and show that
\[
\mu([z_0:z_1])
=
\frac{|z_1|^2}{|z_0|^2+|z_1|^2}
\in[0,1].
\]
Describe the fiber over an interior point and the fibers over the two
endpoints. Write down an explicit representative in the nonnegative
real part over each \(t\in[0,1]\).

\item For \(\PP^2\), take the weights
\[
(0,0),\qquad(1,0),\qquad(0,1).
\]
Write the moment map in homogeneous coordinates. For points in the
interior of the triangle, in the relative interior of an edge, and at
a vertex, determine:
\begin{enumerate}[label=(\alph*),wide=0pt]
\item which homogeneous coordinates are nonzero;
\item the compact stabilizer;
\item the dimension of the moment-map fiber.
\end{enumerate}

\item Regard \(\PP^1\times\PP^1\) as the toric variety of the square
\([0,1]^2\). Prove directly from the product formula that
\[
\mu_{\PP^1\times\PP^1}
=
(\mu_{\PP^1},\mu_{\PP^1}).
\]
Choose a path from the square's interior to a vertex that first meets
an edge. Describe, in order, which circle directions collapse along
the path.

\item Let \(Q\preceq P\) have codimension \(k\). Using the exact
sequence defining \(K_Q\), show that
\[
\dim K_Q=k,
\qquad
\dim(K_N/K_Q)=n-k=\dim Q.
\]
Explain why the dimensions of the moment fibers decrease precisely
when the image crosses into a proper face.

\item Compute the \(h\)-vector in each of the following two ways:
first by orienting edges with a convenient generic height, and then by
using the face-number identity from Section~8.3.
\begin{enumerate}[label=(\alph*),wide=0pt]
\item the \(n\)-simplex;
\item the \(n\)-cube.
\end{enumerate}
Deduce the Betti numbers of \(\PP^n\) and \((\PP^1)^n\).

\item Let \(P\) be a smooth lattice hexagon. Choose a generic height
and prove that its outgoing-edge counts are
\[
2,\quad1,\quad1,\quad1,\quad1,\quad0
\]
in some order. Determine the integral cohomology groups of \(X_P\)
as abelian groups and compute its Euler characteristic. How do these
answers change after one toric blowup at a fixed point?

\item Use the standard fan of \(\PP^n\) to compute \(I_\Sigma\) and
\(J_\Sigma\). Recover
\[
H^*(\PP^n;\ZZ)
\simeq
\frac{\ZZ[H]}{\langle H^{n+1}\rangle},
\qquad
\deg H=2.
\]
Identify geometrically the class \(H^k\).

\item Compute the cohomology ring of \(\PP^1\times\PP^2\) directly
from the product fan. Show that
\[
H^*(\PP^1\times\PP^2;\ZZ)
\simeq
\frac{\ZZ[A,B]}{\langle A^2,B^3\rangle},
\]
and list an additive basis. Verify that its ranks agree with the
\(h\)-vector of the triangular prism.

\item The square and a corner-cut triangle both have \(h\)-vector
\((1,2,1)\). Explain why their associated smooth toric surfaces have
isomorphic cohomology groups as graded abelian groups but nonisomorphic
integral cohomology rings. Formulate the extra information supplied by
the cup product in terms of self-intersection numbers.

\end{exercises}

\clearpage
\section*{Challenge exercise}
\addcontentsline{toc}{section}{Challenge exercise}
\markboth{Challenge exercise}{Challenge exercise}

\noindent
\textbf{Toric surfaces and how much topology remembers.}

\medskip

Let \(X_\Sigma\) be a smooth complete toric surface. Order the rays of
its fan counterclockwise,
\[
\rho_1,\ldots,\rho_s,
\]
write \(u_i\in N\) for the primitive generator of \(\rho_i\), and let
\[
D_i=V(\rho_i),
\qquad
x_i=[D_i]\in H^2(X_\Sigma;\ZZ).
\]
Indices are taken cyclically. Smoothness implies that each consecutive
pair \((u_i,u_{i+1})\) is an oriented lattice basis. Consequently,
for every \(i\) there is a unique integer \(a_i\) such that
\[
\boxed{u_{i-1}+u_{i+1}=a_i u_i.}
\]

The first goal is to recover the complete intersection matrix of the
invariant curves from these relations. The second is to apply the
result to the Hirzebruch surfaces and determine exactly how much of the
parameter \(r\) survives in integral cohomology.

\begin{enumerate}[label=\textbf{Hint \arabic*.},wide=0pt]

\item \textbf{Distinct invariant curves.}

Use the orbit--cone correspondence to show that
\[
D_i\cdot D_j
=
\begin{cases}
1,& \rho_i\text{ and }\rho_j\text{ are adjacent},\\
0,& \rho_i\text{ and }\rho_j\text{ are not adjacent}.
\end{cases}
\]
For adjacent rays, explain why smoothness makes the intersection at the
corresponding fixed point transverse. Translate both cases into
products of the classes \(x_i\) in the Jurkiewicz--Danilov
presentation.

\item \textbf{Self-intersection from three rays.}

Choose \(m_i\in M\) dual to the second vector in the ordered basis
\((u_{i-1},u_i)\), so that
\[
\langle m_i,u_{i-1}\rangle=0,
\qquad
\langle m_i,u_i\rangle=1.
\]
Pair the relation
\[
u_{i-1}+u_{i+1}=a_i u_i
\]
with \(m_i\), and then multiply the principal-divisor relation
\[
\sum_j\langle m_i,u_j\rangle x_j=0
\]
by \(x_i\). Use the nonadjacent Stanley--Reisner relations to make all
but two terms vanish. Prove
\[
\boxed{D_i^2=x_i^2=-a_i.}
\]
Conclude that the cyclic list of primitive ray generators determines
the entire intersection matrix of the invariant curves.

\item \textbf{Four surface checks.}

Apply the formula to the following cyclic fans.
\begin{enumerate}[label=(\alph*),wide=0pt]
\item For \(\PP^2\), use
\[
(1,0),\quad(0,1),\quad(-1,-1)
\]
and show that each invariant line has square \(1\).

\item For \(\PP^1\times\PP^1\), use
\[
(1,0),\quad(0,1),\quad(-1,0),\quad(0,-1)
\]
and show that all four invariant fibers have square \(0\).

\item For \(\Bl_p(\PP^2)\), use
\[
(-1,-1),\quad(1,0),\quad(1,1),\quad(0,1).
\]
Identify \(H=D_{(-1,-1)}\) and \(E=D_{(1,1)}\), and recover
\[
H^2=1,
\qquad
E^2=-1,
\qquad
H\cdot E=0.
\]

\item For the Hirzebruch surface \(\mathbb F_r\), \(r\geq0\), use
\[
u_1=(1,0),\quad
u_2=(0,1),\quad
u_3=(-1,r),\quad
u_4=(0,-1).
\]
Show that
\[
D_1^2=D_3^2=0,
\qquad
D_2^2=-r,
\qquad
D_4^2=r.
\]
\end{enumerate}

\item \textbf{The cohomology ring of \(\mathbb F_r\).}

Introduce variables \(x_1,\ldots,x_4\) for the four rays. Show that
\[
I_\Sigma=\langle x_1x_3,x_2x_4\rangle
\]
and that the character relations are
\[
x_1-x_3=0,
\qquad
x_2+rx_3-x_4=0.
\]
Set
\[
F=x_1=x_3,
\qquad
S=x_2,
\qquad
x_4=S+rF.
\]
Deduce
\[
\boxed{
H^*(\mathbb F_r;\ZZ)
\simeq
\frac{\ZZ[S,F]}{\langle F^2,S^2+rSF\rangle},
\qquad
\deg S=\deg F=2.
}
\]
Normalize the point class by \(S\cdot F=1\), and reconcile the ring
relations with
\[
F^2=0,
\qquad
S^2=-r.
\]

\item \textbf{Changing basis: the even case.}

Suppose \(r=2k\) and define
\[
S'=S+kF.
\]
Compute \((S')^2\) and \(S'F\). Show that the basis \((S',F)\) has
the same intersection form as the two ruling classes on
\(\mathbb F_0\simeq\PP^1\times\PP^1\). Construct an explicit graded
ring isomorphism
\[
H^*(\mathbb F_{2k};\ZZ)
\simeq
H^*(\mathbb F_0;\ZZ).
\]

\item \textbf{Changing basis: the odd case.}

Suppose \(r=2k+1\) and again set \(S'=S+kF\). Show that
\[
(S')^2=-S'F,
\qquad
F^2=0.
\]
Construct an explicit graded ring isomorphism
\[
H^*(\mathbb F_{2k+1};\ZZ)
\simeq
H^*(\mathbb F_1;\ZZ).
\]

\item \textbf{Why parity cannot disappear.}

For \(\alpha=aS+bF\in H^2(\mathbb F_r;\ZZ)\), prove
\[
\alpha^2
=
(-ra^2+2ab)[\mathrm{pt}].
\]
Show that when \(r\) is even, every degree-two integral class has even
self-intersection, whereas when \(r\) is odd, some integral class has
odd self-intersection. Deduce the complete classification
\[
\boxed{
H^*(\mathbb F_r;\ZZ)
\simeq
H^*(\mathbb F_t;\ZZ)
\quad\Longleftrightarrow\quad
r\equiv t\pmod2.
}
\]

\item \textbf{Fan, ring, and polygon.}

On \(\PP^1\times\PP^1\), let \(A\) and \(B\) be the two ruling
classes. Compute \(A\cdot B\) in three ways:
\begin{enumerate}[label=(\alph*),wide=0pt]
\item from adjacency of the corresponding rays;
\item from multiplication in
\(\ZZ[A,B]/\langle A^2,B^2\rangle\);
\item from polygon geometry.
\end{enumerate}
For the last method, take \(P_A\) and \(P_B\) to be unit lattice
segments in the two coordinate directions. Use
\[
P_{A+B}=P_A+P_B
\]
and polarize the identity \(D^2=2\operatorname{Area}(P_D)\) to obtain
\[
D\cdot E
=
\operatorname{Area}(P_D+P_E)
-\operatorname{Area}(P_D)
-\operatorname{Area}(P_E)
\]
for nef divisors in this setting. Verify that the unit square gives
\(A\cdot B=1\).

Finally, explain why the area formula cannot be applied directly to
the exceptional curve \(E\) on \(\Bl_p(\PP^2)\), even though the fan
and cohomology-ring calculations still give \(E^2=-1\).

\end{enumerate}

\begin{threedictionaries}
\begin{center}
\small
\begin{tabularx}{\linewidth}{@{}
  >{\raggedright\arraybackslash}X
  >{\raggedright\arraybackslash}X
  >{\raggedright\arraybackslash}X@{}}
\toprule
\textbf{Fan}
& \textbf{Cohomology ring}
& \textbf{Polygon geometry}
\\
\midrule
adjacent rays
& product is the point class
& transverse boundary meeting
\\
three-ray relation
& self-intersection \(D_i^2=-a_i\)
& change of neighboring edge directions
\\
principal divisors
& integral change of basis
& translate supporting inequalities
\\
Minkowski sum
& polarized cup product
& mixed area
\\
\bottomrule
\end{tabularx}
\end{center}
\end{threedictionaries}

The fans distinguish every \(\mathbb F_r\), but their integral
cohomology rings distinguish only even from odd. Which information in
the parameter \(r\) belongs to the chosen toric structure, which is
visible to algebraic geometry, and which survives in the topology of
the underlying manifold?

\clearpage
\section*{Notes and references}
\addcontentsline{toc}{section}{Notes and references}
\markboth{Notes and references}{Notes and references}

The projective moment map, the nonnegative part of a toric variety,
and the compact-torus quotient are developed in
Cox--Little--Schenck, \S12.2. In particular, Proposition~12.2.3,
Theorem~12.2.5, and Exercises~12.2.8--12.2.9 give the ingredients used
in Sections~8.1 and~8.2.

The orbit filtration, Euler characteristic, vanishing of odd
cohomology, and the \(h\)-vector formula appear in
Cox--Little--Schenck, \S12.3; see especially
Theorems~12.3.9, 12.3.11, and~12.3.12. The attracting-cell language
used here is the toric form of the Bia\l ynicki--Birula decomposition.

Cox--Little--Schenck, Theorem~12.4.4, gives the integral
cohomology-ring presentation for smooth complete toric varieties; this
is the Jurkiewicz--Danilov theorem. The rational version for complete
simplicial toric varieties is Theorem~12.4.1. Our treatment emphasizes
the parallel between the Stanley--Reisner and principal-divisor
relations in cohomology and their earlier roles in the Cox
construction.

The surface calculation in the challenge is a direct application of
that presentation. The identity
\[
u_{i-1}+u_{i+1}=a_i u_i
\quad\Longrightarrow\quad
D_i^2=-a_i
\]
turns a cyclic fan into its intersection matrix. For Hirzebruch
surfaces, an integral change of basis removes all information except
the parity of \(r\) from the cohomology ring. The fan retains the full
integer because it remembers the distinguished torus action and its
invariant boundary, not merely the topology of the underlying
four-manifold.
